\section{Conclusion}
\label{sec:conclusion}

A vulnerability fix is evidence of a convention a project expects its own
code to follow, not just a patch to one function. CVE-2022-4450 shows
this concretely: \texttt{PEM\_read\_bio\_ex} frees two output buffers on
a failure path and leaves the caller-visible pointers dangling, and the
same release-without-reset shape recurs, unrelated to that fix, in four
KDF functions a year earlier. General-purpose static analyzers, run with
their stock rule sets, catch none of the three historical CVEs we
checked.

We presented \tool, which checks a lifecycle rule stored as data against
a real \texttt{tree-sitter} control-flow graph. Protocol~B separates
CVE-2022-4450 from its fix, reproduces the unrelated KDF fix, and flags a
sibling site OpenSSL later fixed. On OpenSSL's current 1{,}155-file tree it flags 10 sites,
of which manual triage confirms 0 new vulnerabilities, 2 potential
recurrences needing reachability review, and 8 harmless deviations.
Protocol~A, for CVE-2023-0215's BIO-chain invariant, reuses the same
parser and CFG and separates its own fix, but its whole-tree run is 0 for
11, for want of ownership tracking: the architecture generalizes across
invariant shapes even where one rule's detection does not. We make no claim beyond this: two rules, one project, with the
threats in Section~\ref{sec:threats}. An ownership/escape analysis for
Protocol A, and a broader historical evaluation, are future work.
